\documentclass[a4paper,10pt]{article}
\usepackage[utf8]{inputenc}
\usepackage[spanish,es-noshorthands]{babel}
\usepackage[margin=1.5cm]{geometry}
\usepackage{tcolorbox}
\usepackage{multicol}
\usepackage{xcolor}
\usepackage{graphicx}
\usepackage{siunitx}
\usepackage{tikz}

\definecolor{primary}{RGB}{33, 64, 154}
\definecolor{secondary}{RGB}{237, 28, 36}
\definecolor{accent}{RGB}{247, 148, 29}
\definecolor{bglight}{RGB}{245, 245, 245}

\tcbuselibrary{skins}

\tcbset{
    infobox/.style={
        enhanced,
        colback=white,
        colframe=primary,
        boxrule=1pt,
        arc=4pt,
        left=5pt,
        right=5pt,
        top=5pt,
        bottom=5pt,
        fonttitle=\bfseries\large,
        coltitle=white,
        attach boxed title to top left={xshift=10pt, yshift=-10pt},
        boxed title style={colback=primary, arc=4pt, boxrule=0pt}
    }
}

\begin{document}

\begin{center}
    \vspace*{0.5cm}
    {\Huge \textbf{\textcolor{primary}{RuView Rescue}}}\\[0.5cm]
    {\LARGE \textbf{Plan de Implementación y Arquitectura}}\\[0.2cm]
    {\large \textit{Grupo de Investigación: Tecnología Venezolana}}\\[1cm]
\end{center}

\begin{tcolorbox}[infobox, title=Descripción del Proyecto, colframe=secondary, boxed title style={colback=secondary}]
Este plan detalla el desarrollo del proyecto \textbf{RuView} aplicado a la búsqueda y rescate en escenarios de desastre (ej. sismos). Su núcleo se basa en usar señales WiFi (CSI) para detectar signos vitales (respiración) a través de escombros, manteniendo estrictamente una \textbf{arquitectura de 3 capas} (Directives $\rightarrow$ Orchestration $\rightarrow$ Execution).
\end{tcolorbox}

\vspace{0.5cm}
\begin{multicols}{2}

\begin{tcolorbox}[infobox, title=Fase 1: Nodo ESP32 (Hardware)]
Desarrollo en C/C++ (ESP-IDF/Arduino) para el ESP32-S3.
\begin{itemize}
    \item \textbf{CSI:} Modo promiscuo para captura de WiFi.
    \item \textbf{Acelerómetro:} MPU6050 vía I2C para registrar ruido mecánico.
    \item \textbf{Transmisión:} UDP hacia el servidor central.
\end{itemize}
\textit{Objetivo:} Capturar datos físicos crudos con ultra-baja latencia.
\end{tcolorbox}

\begin{tcolorbox}[infobox, title=Fase 2: Motor Central (3 Capas)]
Servidor local (Raspberry Pi/PC).
\begin{itemize}
    \item \textbf{Capa 1 (Directives):} \texttt{ruview\_rescue.yaml}. Parámetros, umbrales y puertos.
    \item \textbf{Capa 2 (Orchestration):} Script Python UDP que enruta la data y despacha alertas.
    \item \textbf{Capa 3 (Execution):} Script Python para filtrado matemático. Uso de filtro adaptativo NLMS para restar vibración de maquinaria pesada y FFT para detectar respiración (0.2 - 0.5 Hz).
\end{itemize}
\end{tcolorbox}

\begin{tcolorbox}[infobox, title=Fase 3: Diseño KiCAD]
Diseño de la Placa de Circuito Impreso (PCB) para el nodo táctico.
\begin{itemize}
    \item Aislamiento contra ruido de fuentes conmutadas.
    \item Protección del acelerómetro.
    \item Optimización con antenas direccionales.
\end{itemize}
\end{tcolorbox}

\begin{tcolorbox}[infobox, title=Fase 4: Inteligencia Artificial]
\begin{enumerate}
    \item \textbf{Data:} Recolección masiva con v1 (Python).
    \item \textbf{Nube:} Entrenamiento de Redes Neuronales (PyTorch/TensorFlow).
    \item \textbf{Cuantización:} Exportación a modelo ultraligero (TFLite).
    \item \textbf{Edge AI:} Inferencia local en Raspberry Pi migrando el motor de ejecución a \textbf{Rust (v2)} para máxima eficiencia (0\% dependencia de internet).
\end{enumerate}
\end{tcolorbox}

\end{multicols}

\vspace{1cm}
\begin{center}
    \textit{Arquitectura de Monitoreo Espacial y Procesamiento de Señales Avanzadas.}
\end{center}

\end{document}
